\section{A computed boundary of the pencil incidence space}
\label{sec:computed-boundary-v158}
We give the incidence compactification a precise universal property and
compute a family of its boundary fibres in every dimension. The new
boundary is selected by one power base ideal, not by pulling an
unrelated algebra back to a parameter space.

\begin{proposition}[The normal main-component flattening property]
\label{prop:flattening-property-v158}
Let $G^\circ\subset G=\Gr(2,\Sym^2V)$ and
$\widehat G\to G$ be as in Proposition~\ref{prop:pencil-compactification-v157}.
Among normal integral schemes $T$ with a dominant morphism $f:T\to G$
for which the closure of the lifted pencil family over $f^{-1}(G^\circ)$
is an embedded flat family in $\mathrm{CQ}(V)\times T$ with Hilbert
polynomial $\binom n2 m+1$, there is a unique morphism $T\to\widehat G$
over $G$, and that family is the pullback of its universal family.
Consequently $\widehat G$ is the minimal normal dominant modification
with this specified embedded-flat-closure property.
\end{proposition}
\begin{proof}
The flat embedded family gives a morphism to the indicated Hilbert
scheme. Together with $f$ it agrees, over a dense open of $T$, with
the graph defining the reduced graph closure. Since $T$ is integral,
its morphism factors through this closed graph scheme: each equation
vanishing on the dense open vanishes on $T$. The induced morphism is
dominant. A dominant morphism from a normal integral scheme to an
integral scheme factors uniquely through the normalization: locally,
an element integral over the target ring becomes integral over the
source ring and lies in its fraction field, hence belongs to the
normal source ring. Thus the morphism factors uniquely through
$\widehat G$. The Hilbert-scheme representing property identifies
the families. Conversely its pulled-back universal family is flat;
flatness over an integral base makes it the schematic closure of its
restriction to the dense open. Uniqueness follows also from separatedness.
\end{proof}

The hypothesis of dominance is part of this property; this proposition
does not assert that normalization commutes with arbitrary base change.
The property is a standard consequence of the Hilbert functor and
normalization. Its usefulness here is that the following calculation
identifies a nontrivial fibre of the resulting modification.

Assume $n\geq3$. Choose distinct nonzero numbers $a_3,\ldots,a_n$ and put
$D=\Spec\C[[\tau]]$. On $\Pj^1_D$, consider the symmetric pencil
\begin{equation}\label{eq:collision-family-v158}
 A_\tau(s,t)=\diag\bigl(s,s+\tau t,
                           t-a_3s,\ldots,t-a_ns\bigr).
\end{equation}
For $n=3$ there is just one number $a_3$. Write $X=\Pj^1_D$ and let
$o$ be the point $s=\tau=0$ in the special fibre. Its generic pencil
has $n$ distinct rank-$(n-1)$ spectral points, so defines a point of
$G^\circ$.

\begin{theorem}[A transverse corank-two tail and its power maps]
\label{thm:hilbert-tail-v158}
The closure of the generic lifted curve of
\eqref{eq:collision-family-v158} in $\mathrm{CQ}(V)\times D$ is
\begin{equation}\label{eq:tail-blowup-v158}
                         Z=\operatorname{Bl}_o X.
\end{equation}
It is flat over $D$, and its special fibre is the reduced nodal union
$C\cup F$, where $C\simeq\Pj^1$ is the lifted special pencil and
$F\simeq\Pj^1$ is an exceptional tail. There are no embedded components.
If $H_q$ is the exterior-power hyperplane bundle of order $q$, then
\begin{equation}\label{eq:tail-degrees-v158}
 (\deg_C H_q,\deg_F H_q)=
 \begin{cases}(q,0),&1\leq q\leq n-2,\\
              (n-2,1),&q=n-1.
 \end{cases}
\end{equation}
Thus the special fibre has the same Hilbert polynomial
$\binom n2m+1$ as the generic lifted line.

The tail lies in $E_{n-2}$, its two intersections with $E_{n-1}$
are distinct and reduced, and its attachment to $C$ is neither of
those two points. For any $h\geq1$, let $L_p$ be the hyperplane
bundle of the resolved $p$th power map. Then
\begin{equation}\label{eq:tail-power-degree-v158}
 \deg_F L_p=
 \begin{cases}
 0,&p\leq h(n-2),\\
 \min\{b,2h-b\},&p=h(n-2)+b,\quad 0\leq b\leq2h.
 \end{cases}
\end{equation}
When this degree is positive, the restricted map on $F$ is the
complete Veronese map of that degree into its linear span.
Most directly, in the local coordinate $z=s/t$ at $o$,
\begin{equation}\label{eq:one-power-centre-v158}
 J_{h(n-2)+1}(A_\tau)=(z,\tau).
\end{equation}
Hence a single multiplication power ideal selects the blow-up centre
and its exceptional component.
\end{theorem}
\begin{proof}
The rank is at least $n-2$ on $X$. In a neighbourhood of $o$ the last
$n-2$ diagonal entries are units; all other spectral points have rank
$n-1$ and are pairwise distinct. Consequently the extended minor ideals
$I_q$ are the unit ideal for $q\leq n-2$, while
\[
 I_{n-1}=(z,z+\tau)=(z,\tau),\qquad
 I_n=(u z(z+\tau))
\]
near $o$, with $u$ a unit. Away from $o$, $I_{n-1}$ is a unit ideal;
$I_n$ is everywhere invertible. The simultaneous exterior-power graph
on the integral surface $X$ is therefore the blow-up of $(z,\tau)$.
This is an application of the graph-by-generators lemma, not a claim
that arbitrary base change of a blow-up is a blow-up. Since the
universal family of pencil lines embeds $X$ into $\Pj(\Sym^2V)\times D$,
this graph embeds into $\mathrm{CQ}(V)\times D$ and is the required
scheme-theoretic closure.

On the chart with $\tau=zv$, the special fibre has equation $zv=0$;
on the other chart, $z=\tau w$, it has equation $\tau=0$.
These charts prove flatness over the DVR, multiplicity one of both
components, their transverse intersection, and absence of embedded
components. The map on $C$ is the unique lift of the special pencil.
The $q$-minors of that pencil have degree $q$. They have no common
factor for $q\leq n-2$, and precisely one factor $s$ for $q=n-1$.
This proves their degrees on $C$. On $F$, all lower maps are constant,
and the $(n-1)$st exterior map is linear. This gives
\eqref{eq:tail-degrees-v158}. For $L=\bigotimes H_q$, the degrees
on $C$ and $F$ add to $\binom n2$. The normalization sequence of
$C\cup F$ subtracts one point, so its Euler characteristic is one
and its Hilbert polynomial is as stated. Thus it is the Hilbert
limit used by $\widehat G$; equivalently the generic morphism from
$D$ extends to $\widehat G$ by properness, with this universal fibre.

Use homogeneous coordinates $[\alpha:\beta]=[z:\tau]$ on $F$.
After removing the exceptional factor, the residual quadratic form
on the two-dimensional kernel is
$\diag(\alpha,\alpha+\beta)$. This is a line in
$\mathrm{CQ}(\C^2)=\Pj(\Sym^2\C^2)$, and gives the claimed image
of $F$ in the fibre of $E_{n-2}$. Its determinant $\alpha(\alpha+\beta)$ has two
simple zeros $[0:1]$ and $[-1:1]$, which are its intersections with
$E_{n-1}$. The main component attaches at $[1:0]$, where this residual
form is nondegenerate.

For clarity, the whole list of extended power ideals near $o$ is
\begin{equation}\label{eq:tail-all-ideals-v158}
 J_p=
 \begin{cases}
 \OO,&p\leq h(n-2),\\
 (z,\tau)^b,&p=h(n-2)+b,\ 0\leq b\leq h,\\
 (z,\tau)^{2h-b}[z(z+\tau)]^{b-h},
       &p=h(n-2)+b,\ h\leq b\leq2h,
 \end{cases}
\end{equation}
up to an invertible factor. These are actual ideals, by
Theorem~\ref{thm:universal-power-v157}. In particular $b=1$ gives
\eqref{eq:one-power-centre-v158} for every $h$.

It remains to check the linear system on the tail, not only its
degree. On diagonal matrices the coefficient polynomials of a $p$th
power span exactly the monomials of total degree $p$ in the diagonal
entries with each exponent at most $h$. For containment, use the
torus-weight bound in Lemma~\ref{lem:bounded-row-v157}; for the reverse
inclusion, take the coefficients of the independent diagonal entries
of $X$ in $\det(X+uA)^h$. Thus, for $p=h(n-2)+b$, the leading normal
coefficients use exponent $h$ in each of the $n-2$ invertible entries
and span
\[
  \alpha^j(\alpha+\beta)^{b-j},\qquad
           \max\{0,b-h\}\leq j\leq\min\{b,h\}.
\]
If $b\leq h$ these span every
binary form of degree $b$. If $b>h$, remove the common factor
$[\alpha(\alpha+\beta)]^{b-h}$; the remaining monomials span every binary form of
degree $2h-b$. Lower powers have nonzero constant normal term and
are constant on $F$. This proves the complete-linear-system assertion
and \eqref{eq:tail-power-degree-v158}.
\end{proof}

\begin{remark}[What the failure family contributes]
\label{rem:tail-integration-v158}
Along \eqref{eq:collision-family-v158}, the original finite failure
algebras form the finite locally free family already constructed in
Corollary~\ref{cor:finite-flat-family}. Its pullback to $F$ is constant;
we do not count that pullback as a new interaction. Instead, its first
relation canonically recovers the source and the pencil line, and the
source-normalized envelope has the total-base ideal
\eqref{eq:one-power-centre-v158}. Resolving that multiplication ideal
creates $F$, and the other products prescribe all the linear systems
\eqref{eq:tail-power-degree-v158} on it. A degeneration arc, not only
its closed failure algebra, supplies this total-base ideal. Thus the
calculation distinguishes information in the family from information
in the special point and supplies an explicit modular Hilbert limit.
\end{remark}
